\section{Method~I: oscillatory residual-stress realisation}
\label{sec:methodI}

The goal of this section is to construct a divergence-free field $u=u_B+w+v$ whose residual extends smoothly through $t=1$, while $u$ remains comparable to $u^{(0)}$ on a collapsing ring. Here $u_B$ is an axisymmetric background correcting $u^{(0)}$ to every order in the core, $w$ is a superposition of oscillatory pulses, and $v$ is a rapidly convergent remainder.

\subsection{Background correctors}
\label{sec:background}

\begin{construction}[Background expansion]
\label{con:background}
Starting from $u^{(0)}$, we construct axisymmetric correctors $u^{(n)}$, $n\ge 1$, so that
\[
u_B:=\sum_{n=0}^N u^{(n)},
\qquad
p_B:=\sum_{n=0}^N p^{(n)},
\]
for a finite $N=N(h)$ large enough that the leftover core residual is $O(q^{M})$ for an arbitrarily large $M$. Each $u^{(n)}$ is of size $O(q^{-A+2nh})$ in the core, corresponding to the expansion parameter $q^{2h}$ of Remark~\ref{rem:q2h}.
\end{construction}

The linearised operator about $u^{(0)}$ in similarity variables is a degenerate elliptic/parabolic operator $\mathcal L$ on the rectangle $[0,X_c]\times[-\eta_c,\eta_c]$, with the axis regularity of Appendix~\ref{app:axis} imposed as boundary conditions at $X=0$. At each order one solves
\[
\mathcal L\bigl(u^{(n)},p^{(n)}\bigr)
= -\mathcal N_{<n}-\mathcal A_{\mathrm{ax}}u^{(n-1)},
\]
where $\mathcal N_{<n}$ collects quadratic interactions of lower-order terms and $\mathcal A_{\mathrm{ax}}$ is axial diffusion (the term omitted from the leading balance). Solvability follows from the same coercivity that produces the leading profile, up to the finite-dimensional kernel generated by scaling and by vertical translation, both of which are fixed by our choice of singular point and of $A,D$.

\begin{proposition}[Corrected background]
\label{prop:background}
There exists a finite expansion $u_B,p_B$ as in Construction~\ref{con:background}, smooth and compactly supported in similarity space on $|\eta|\le\eta_c$, $X\le X_b$, such that
\begin{equation}
\label{eq:RB}
\Rmom{u_B}{p_B}
=-\divv T_{\ann}
+\mathcal E_{\infty},
\end{equation}
where $T_{\ann}$ is supported in the annulus, takes values in the cone $\mathcal K$, satisfies~\eqref{eq:Tann-size}, and $\mathcal E_{\infty}$ extends smoothly through $t=1$. Moreover $u_B-u^{(0)}=O(q^{-A+2h})$ in $L^\infty(\mathcal C_\tau)$, so the blowup ring~\eqref{eq:blowup-ring} persists:
\begin{equation}
\label{eq:uB-blowup}
\|u_B(\cdot,1-\tau)\|_{L^\infty(\mathcal C_\tau)}\ \ge\ \tfrac12 E(X_*,0)\,\tau^{-1/2-h}
\qquad\text{for all sufficiently small }\tau.
\end{equation}
\end{proposition}

\begin{proof}
The expansion is finite, each corrector is smooth on the similarity rectangle, and the leftover after $N$ steps is $O(q^{-A-1+2Nh})$. Choosing $2Nh>A+1+M$ makes this $O(q^{M})$ for any prescribed $M$; multiplying by a cutoff in $(X,\eta)$ produces a field vanishing to order $M$ at the origin at time $1$. Derivatives cost negative powers of $q$, so taking $M$ larger than any fixed derivative order yields a $C^\infty$ extension by $0$. The annular identity is Proposition~\ref{prop:Tann} applied to $u_B$ instead of $u^{(0)}$; the stress changes by $O(q^{-2A+2h})$, which remains inside the open cone $\mathcal K$ for $h$ small and $q$ small. The lower bound~\eqref{eq:uB-blowup} is~\eqref{eq:blowup-ring} minus the first corrector.
\end{proof}

\subsection{Exact increment identity}
\label{sec:increment}

If one adds a divergence-free increment $w$ with pressure increment $\pi$, the residual changes by an exact identity that isolates the {R}eynolds stress.

\begin{lemma}[Increment identity]
\label{lem:increment}
Let $u,p$ be smooth and $\nabla\cdot u=0$, and let $w,\pi$ be smooth with $\nabla\cdot w=0$. Then
\begin{equation}
\label{eq:increment}
\Rmom{u+w}{p+\pi}
=\Rmom{u}{p}
+\partial_t w+(u\cdot\nabla)w+(w\cdot\nabla)u-\Delta w+\nabla\pi
+\divv(w\otimes w).
\end{equation}
In particular, writing $\overline{w\otimes w}$ for an angular (or phase) average and $w\otimes w-\overline{w\otimes w}$ for the fluctuation,
\[
\divv(w\otimes w)=\divv\overline{w\otimes w}+\divv\bigl(w\otimes w-\overline{w\otimes w}\bigr).
\]
\end{lemma}

\begin{proof}
Expand $(u+w)\cdot\nabla(u+w)$ and $\Delta(u+w)$. The identity $\divv(w\otimes w)=(w\cdot\nabla)w$ holds because $\nabla\cdot w=0$.
\end{proof}

The strategy is to choose $w$ so that $\divv\overline{w\otimes w}=+\divv T_{\ann}$ to leading order, while $w$ itself solves a linearised evolution along $u_B$ to high accuracy, so that the linear increment
\[
L_{u_B}(w,\pi):=\partial_t w+(u_B\cdot\nabla)w+(w\cdot\nabla)u_B-\Delta w+\nabla\pi
\]
is smaller than the mean stress, or is cancelled by a later corrector. The fluctuation $\divv(w\otimes w-\overline{w\otimes w})$ is high frequency and is removed by a {C}alder{\'o}n--{Z}ygmund corrector, as in convex integration, except that here the amplitude is not arbitrary: it is slaved to amplification by the background shear.

\subsection{Two pulse families}
\label{sec:pulses}

\begin{construction}[Ring pulses]
\label{con:pulses}
A \emph{pulse} of family $\sigma\in\{+1,-1\}$ is a divergence-free field of the form
\begin{equation}
\label{eq:pulse}
w_{\sigma,j}(x,t)
=\Re\Bigl(
a_{\sigma,j}(r,z,t)\,
\Psi_\sigma\bigl(\phi_{\sigma,j}(r,z,t),\theta\bigr)
\,e^{i k_{\sigma,j}\vartheta_{\sigma,j}(r,z,t)}
\Bigr),
\end{equation}
supported in a space-time region that, at each frozen $t$, is a complete solid ring: an interval in $r$, an interval in $z$, and the full circle in $\theta$. The index $j\in\N$ labels dyadic generations as $\tau\downarrow 0$. Here:
\begin{itemize}[leftmargin=1.2em]
\item $\phi_{\sigma,j}$ is a slowly varying envelope, equal to $1$ on a slightly smaller ring and vanishing near the ring's boundary;
\item $\vartheta_{\sigma,j}$ is a real phase whose gradient (the local wavevector) is essentially radial-axial, of length $\lambda_{\sigma,j}^{-1}$ with $\lambda_{\sigma,j}\ll \ell_r$;
\item $\Psi_\sigma$ is a vector polarisation, divergence-free to leading WKB order, with vanishing $\theta$-mean;
\item $a_{\sigma,j}$ is a complex amplitude.
\end{itemize}
The two families are distinguished by the ratio of their mean radial angular-momentum flux to their mean radial axial-momentum flux. This is the two-dimensional stress cone of Appendix~\ref{app:cone}.
\end{construction}

\paragraph{Seeding, amplification, and viscous cutoff.}
Each pulse is seeded by an external force of size $\exp(-c/\tau^{\gamma})$ for a small $\gamma>0$, hence smaller than every power of $\tau$ and invisible to the $C^\infty$ extension of the residual. The background shear then amplifies the pulse. A schematic of the mechanism, sufficient for the scaling, is as follows.

In a rotating flow whose angular velocity $\Omega=u_\theta/r$ decreases rapidly with $r$, an azimuthal surplus produces a centrifugal imbalance and a radial velocity; that radial displacement, in a region with $\partial_r\Omega<0$, increases the surplus relative to the surroundings. This is the centrifugal instability of columnar vortices~\cite{leibovichstewartson1978}. Axial shear $\partial_r u_z$ supplies a second, independent amplification mechanism, needed near $z=0$ where the axial bias of \ref{P4} has moved the zero of $u_z$ off the plane of weakest rotational amplification.

At the same time the differential rotation winds the wavevector, shortening the radial wavelength. Viscous damping grows like $|\xi|^2\sim \lambda^{-2}$. We choose the initial wavelength of generation $j$, placed at time $\tau_j=2^{-j}$, so that amplification dominates on a short time interval of length $\asymp\tau_j$, after which damping wins and the pulse is cut off with tails smaller than $\tau_j^{N}$ for every $N$.

\begin{lemma}[WKB amplitude, schematic form]
\label{lem:wkb}
Along particle trajectories of $u_B$, the leading amplitude of a pulse with local wavevector $\xi$ satisfies
\begin{equation}
\label{eq:wkb}
\Dt \log |a|
=\mathcal G(u_B,\hat\xi)
-|\xi|^2
+O\bigl(\lambda\,|\nabla u_B|\bigr),
\end{equation}
where $\mathcal G$ is a homogeneous function of degree one in $\nabla u_B$, given by the growth rate of the {L}ifschitz--{H}ameiri / {F}riedlander--{V}ishik geometric-optics problem~\cite{lifschitz1991,friedlandervishik1991}, and $\hat\xi=\xi/|\xi|$. On the annular support, \ref{P6} yields $\mathcal G\ge \kappa q^{-1}$ for a geometric constant $\kappa>0$ depending only on the profile, while $|\xi|^2$ is initially $O(q^{-1+\delta})$ and finally $O(q^{-1-\delta})$ after winding, for a small $\delta>0$.
\end{lemma}

The proof of~\eqref{eq:wkb} is the standard WKB substitution of~\eqref{eq:pulse} into the linearised operator $L_{u_B}(w,\pi)=0$, collecting terms of order $|\xi|$. We treat~\eqref{eq:wkb} as part of the analytic content of Assumption~\ref{ass:pulses} below; the algebraic form is recorded because it makes the construction of generations canonical.

\paragraph{Mean stress.}
Because $\Psi_\sigma$ has zero $\theta$-mean, $\overline{w_{\sigma,j}}=0$. The quadratic products need not average to zero. For a single family one obtains a rank-one contribution to $\overline{w\otimes w}$ whose direction in the $(r\theta, rz)$ flux plane is determined by $\sigma$. Two families with linearly independent flux ratios span a cone. Realising a prescribed annular stress is then a question of choosing local intensities $|a_{\sigma,j}|^2$.

\begin{proposition}[Leading mean stress]
\label{prop:mean-stress}
Let $T_{\ann}$ take values in the interior of $\mathcal K$ as in Appendix~\ref{app:cone}. There exist intensity functions $\iota_{\pm}\ge 0$, smooth on the annular support, such that if two pulse families are superimposed with $|a_{\pm,j}|^2=\iota_{\pm}\,q^{-2A}$ on their active set at generation $j$, then
\begin{equation}
\label{eq:mean-stress}
\divv\overline{w_j\otimes w_j}
=\divv T_{\ann}
+O\bigl(q^{-2A}\lambda_j/\ell_r\bigr)
+O\bigl(q^{-2A}\,2^{-j}\bigr)
\end{equation}
on the space-time support of generation $j$, and the left-hand side vanishes off that support. The error $O(q^{-2A}\lambda_j/\ell_r)$ is the WKB remainder; the error $O(q^{-2A}2^{-j})$ is the discrepancy between $T_{\ann}(\cdot,t)$ and its value at the centre of the generation.
\end{proposition}

\begin{proof}[Proof of the algebraic part]
Pointwise, $\overline{w\otimes w}$ is a linear combination
\[
\iota_+ S_+(\hat\xi_+)+\iota_- S_-(\hat\xi_-)
\]
of two symmetric tensors determined by the polarisations. Appendix~\ref{app:cone} shows that $\mathcal K$ is the cone generated by $S_\pm$ as the wavevector varies in an open set of directions compatible with~\ref{P6}. Choosing $\iota_\pm$ is therefore a $C^\infty$ linear algebra problem on the annulus, with a smooth positive solution. Multiplying by the slowly varying envelope and summing over a partition of unity in time with scale $\tau_j$ produces~\eqref{eq:mean-stress} up to the two errors indicated, which are estimated in Assumption~\ref{ass:pulses}.
\end{proof}

\paragraph{Separation of supports.}
Pulses of a given generation occupy disjoint rings, and distinct generations occupy disjoint time intervals except for rapidly decaying tails. An auxiliary large torus, on which the construction is first performed before being mapped back to a small ball, makes the disjointness geometric rather than analytic; this is the device of~\cite{openai2026ns}, and we adopt it at the level of the scheme.

\subsection{The remaining analytic assumption for Method~I}

Everything in Method~I that is not already proved is collected here. It is a finite list of estimates, not an open-ended request.

\begin{assumption}[Pulse iteration]
\label{ass:pulses}
There exist pulse generations $(w_j,\pi_j)_{j\ge J_0}$ as in Construction~\ref{con:pulses}, and correctors $(v_m,\varrho_m)_{m\ge 0}$, such that the following hold with implied constants independent of $j$ for all sufficiently large $j$.
\begin{enumerate}[label=\textup{(A\arabic*)}]
\item \label{A1} \emph{Linear evolution.} $\|L_{u_B}(w_j,\pi_j)\|_{C^N}\le C_{N,M}\tau_j^{M}$ for every $N,M$, after including the exponentially small seed force in the residual. Equivalently, the WKB remainder and the cut-off tails are super-smooth at $t=1$.
\item \label{A2} \emph{Mean stress.} Identity~\eqref{eq:mean-stress} holds, and $\sum_j\bigl(\overline{w_j\otimes w_j}-T_{\ann}\chi_j\bigr)$ extends smoothly through $t=1$, where $\chi_j$ is the generation cutoff.
\item \label{A3} \emph{Fluctuation.} The high-frequency fluxes $\divv(w_j\otimes w_j-\overline{w_j\otimes w_j})$ are the divergence of a stress whose $C^N$ norms on the complement of a $q$-neighbourhood of the origin are $O(\tau_j^{M})$ for every $N,M$, and which is cancelled in that neighbourhood by the first corrector $v_0$ up to a smoother error.
\item \label{A4} \emph{Amplitude versus background.} $\sum_j\|w_j\|_{L^\infty(\mathcal C_\tau)}=O(\tau^{-A+h/2})$, so that $w$ cannot cancel the blowup ring~\eqref{eq:uB-blowup}. Off the annulus, $w_j=0$.
\item \label{A5} \emph{Correctors.} There is a sequence $v_m$ with $\|v_m\|_{C^N}=O(\tau^{m\beta-C_N})$ for some $\beta>0$, such that $v:=\sum_m v_m$ converges in $\Cinf$ on $\R^3\times[0,1)$ and
    \[
    \Rmom{u_B+w+v}{p_B+\pi+\varrho}
    \quad\text{extends to a function in }\Ccinf(\R^3\times(0,1];\R^3)
    \]
    before spatial cutoff. Here $w=\sum_j w_j$.
\item \label{A6} \emph{Incompressibility and support.} Each $w_j$ and $v_m$ is divergence-free, and supported in a fixed compact set of $\R^3$ independent of $j,m$ (after the localisation of Section~\ref{sec:cutoff}).
\end{enumerate}
\end{assumption}

\begin{remark}
\label{rem:assI-status}
Estimates \ref{A1}--\ref{A6} are precisely the content of Sections~6--9 of~\cite{openai2026ns}, rewritten so as to apply to any background satisfying Proposition~\ref{prop:background}. We do not reproduce those pages in this draft. The reduction above makes the dependence on them completely explicit: if \ref{A1}--\ref{A6} hold, Theorem~\ref{thm:methodI} is proved by the localisation of Section~\ref{sec:cutoff} and by Lemmas~\ref{lem:residual-criterion}--\ref{lem:periodise}.
\end{remark}

\subsection{Summation before cutoff}

\begin{proposition}[Pre-localised Method~I field]
\label{prop:prelocal-I}
Grant Assumption~\ref{ass:pulses}. Set $u^{\mathrm I,\#}:=u_B+w+v$ and let $p^{\mathrm I,\#}$ be the corresponding pressure. Then $u^{\mathrm I,\#}$ is smooth and divergence-free on $\R^3\times[0,1)$,
\[
\limsup_{t\uparrow 1}\|u^{\mathrm I,\#}(t)\|_{L^\infty}=\infty,
\qquad
\sup_{t\in[0,1)}\|u^{\mathrm I,\#}(t)\|_{L^2}<\infty,
\]
and $\Rmom{u^{\mathrm I,\#}}{p^{\mathrm I,\#}}$ extends smoothly through $t=1$ as a compactly supported function of $(x,t)$ in a neighbourhood of the origin (not yet compactly supported in a $\tau$-independent ball, nor vanishing near $t=0$).
\end{proposition}

\begin{proof}
Smoothness and divergence-free are \ref{A5}--\ref{A6}. The $L^\infty$ blowup is~\eqref{eq:uB-blowup} plus \ref{A4}. The $L^2$ bound on the core is Proposition~\ref{prop:scales}; the pulses live in a set of volume $O(\tau^{3/2-h})$ with amplitude at most $O(\tau^{-A+h/2})$, hence contribute $O(\tau^{1/2-2h})$ to the energy, which is bounded; the heat exterior is smooth through $t=1$ at each $r>0$ and will be cut off in the next subsection, so its energy is bounded as well. The residual extends by \ref{A5} and Proposition~\ref{prop:background}.
\end{proof}

\subsection{Spatial cutoff, rest initial datum, and the force}
\label{sec:cutoff}

It remains to: (i)~cut the field off outside a ball, without destroying smoothness of the residual through $t=1$; (ii)~arrange $u(\cdot,0)=0$; (iii)~rescale viscosity and periodise.

\begin{lemma}[Spatial cutoff]
\label{lem:cutoff-space}
Let $\psi\in\Ccinf(\R^3)$ equal $1$ on the unit ball and vanish outside the ball of radius $2$. For $\rho>0$ small, set
\[
u_\rho(x,t):=\psi(x/\rho)\,u^{\mathrm I,\#}(x,t)+Z_\rho(x,t),
\]
where $Z_\rho$ is a {B}ogovskii correction, compactly supported in $\{1\le |x|/\rho\le 2\}$, restoring $\nabla\cdot u_\rho=0$ (possible because the flux of $\psi(x/\rho)u^{\mathrm I,\#}$ through spheres is zero by $\nabla\cdot u^{\mathrm I,\#}=0$ and compact support in the angular variables). Then for $\rho$ small enough that the core $\mathcal C_\tau$ lies in $\{|x|<\rho/2\}$ for all small $\tau$, the inner blowup is unchanged, and $\Rmom{u_\rho}{p_\rho}$ differs from $\psi(x/\rho)\Rmom{u^{\mathrm I,\#}}{p^{\mathrm I,\#}}$ by terms supported in the annulus $\{\rho\le |x|\le 2\rho\}$. On that annulus, $u^{\mathrm I,\#}$ and all its derivatives have smooth limits as $t\uparrow 1$, because the only singularity is at the origin. Hence the additional residual extends smoothly through $t=1$.
\end{lemma}

\begin{proof}
The only possible obstruction would be a singularity of $u^{\mathrm I,\#}$ at some $x\neq 0$, $t=1$. In the heat exterior, Proposition~\ref{prop:heat} and the matching of Appendix~\ref{app:moments} give smoothness at each $r>0$. In the annulus and core, $q\gtrsim |x|^{C}$ away from the origin, so $q$ is bounded below on $\{|x|\ge\rho\}$ uniformly as $\tau\downarrow 0$. Thus all negative powers of $q$ remain bounded there.
\end{proof}

\begin{lemma}[Turn-on from rest]
\label{lem:turnon}
Let $\eta\in\Cinf(\R)$ with $\eta(t)=0$ for $t\le 0$ and $\eta(t)=1$ for $t\ge t_*$, for a small $t_*>0$. Replacing $(u_\rho,p_\rho)$ on $[0,t_*]$ by the unique smooth {N}avier--{S}tokes solution with force $\eta(t) f_\rho$ and initial datum $0$, where $f_\rho$ denotes the already constructed residual for $t>t_*$, yields a global-in-time-up-to-$1$ field starting from rest. For $t_*$ small the inner construction on $[t_*,1)$ is untouched.
\end{lemma}

\begin{proof}
Short-time existence from rest with a smooth compactly supported force is classical~\cite{kato1972}. Choose $t_*$ smaller than the first time at which any pulse generation of Assumption~\ref{ass:pulses} is seeded, and smaller than the time at which the core has shrunk inside the cutoff ball. Concatenation at $t_*$ is smooth because both pieces solve the same equation with the same force.
\end{proof}

An equivalent and slightly cleaner procedure, which we adopt, is to run the whole inner construction on the time interval $[0,1)$ with a time cutoff that vanishes near $t=0$ inserted into the seed forces and into $T_{\ann}$. Then $u\equiv 0$ near $t=0$ automatically.

\begin{proof}[Proof of Theorem~\ref{thm:methodI}]
Grant Assumption~\ref{ass:pulses}. Proposition~\ref{prop:prelocal-I} and Lemmas~\ref{lem:cutoff-space}--\ref{lem:turnon} produce, at viscosity $1$, a field satisfying all hypotheses of Lemma~\ref{lem:residual-criterion}. Lemma~\ref{lem:energy} gives the energy bound. Lemma~\ref{lem:rescale} gives every $\nu>0$. A preliminary homothety $x\mapsto \lambda x$ with $\lambda$ large, followed by an affine transformation of time restoring $T_*=1$, places the support in $(-1/4,1/4)^3$. Corollary~\ref{cor:torus} gives alternative (D).
\end{proof}

This completes Method~I, conditionally on Assumption~\ref{ass:pulses}.
